\documentclass[12pt]{amsart}

\usepackage[margin=1in]{geometry}
\usepackage{amsfonts}
\usepackage{amssymb}
\usepackage{hyperref}
\hypersetup{
    colorlinks=true,
    linkcolor=black,
    filecolor=green,
    urlcolor=cyan,
    citecolor=black,
    pdftitle={Service Statement},
}
\urlstyle{same}


\begin{document}

\begin{center}
  {\large\textsc{Service Statement}}\\[5pt]
  \small\textsc{Edward A. Varvak, University of Chicago}
\end{center}

Most of my service so far has been in the mathematics department at the
University of Illinois Chicago (UIC), where I began my Ph.D. There, I have
tried to give graduate students more places to learn mathematics from one
another, and more people to turn to for advice.

For two semesters, from Fall 2025 through Spring 2026, I co-organized the
Graduate Algebraic Geometry Seminar, a student-run seminar in which graduate
students give talks on algebraic geometry. I also organized two learning
seminars for graduate students: a reading group on the Weil conjectures in
Fall 2025, and a seminar on Lie groups in Fall 2024. For each, I chose the
topic, broke it into a dozen talks, assigned them to participants, and
helped speakers prepare when they needed it.

I have also been a panelist at two events in the UIC Department of
Mathematics, Statistics, and Computer Science: a geometry advising panel,
where first- and second-year graduate students heard how to plan a Ph.D.\ in
geometry, from passing preliminary exams to choosing an advisor and starting
research; and a visit day panel, where I answered prospective students'
questions about the program and graduate life. More informally, I have been a
mentor to younger graduate students as they navigate graduate school.

Much of my teaching has been in the same community. I spent five semesters
as a teaching assistant at UIC, a Hispanic-Serving Institution where many
students are the first in their families to attend college. Outside the
university, I judged QED, Chicago's Youth Math Symposium, run by the Math
Circles of Chicago, at which students in grades 5--12 present original
projects in mathematics and computer science.

At Harvard, I plan to organize a learning seminar on Lagrangian fibrations
and Hitchin systems, open to graduate students and postdocs across the
Boston area, and to post its notes publicly so that students without a local
expert can follow along. I would also like to continue outreach by
volunteering with local math circles and competitions such as HMMT.

\end{document}
